\documentclass[11pt,a4paper]{article}
% lesson-preamble.tex -- Shared preamble for all Phase 00 lessons
% Usage: % lesson-preamble.tex -- Shared preamble for all Phase 00 lessons
% Usage: % lesson-preamble.tex -- Shared preamble for all Phase 00 lessons
% Usage: \input{../lesson-preamble} (from subdomain directories)

% ---- Encoding and fonts ----
\usepackage[utf8]{inputenc}
\usepackage[T1]{fontenc}

% ---- Mathematics ----
\usepackage{amsmath,amssymb,amsthm}
\usepackage{mathtools}

% ---- Page geometry ----
\usepackage[margin=1in,headheight=14pt]{geometry}

% ---- Hyperlinks ----
\usepackage[colorlinks=true,linkcolor=red,citecolor=blue,urlcolor=blue]{hyperref}

% ---- Lists ----
\usepackage{enumitem}

% ---- Colors and boxes ----
\usepackage{xcolor}
\usepackage[theorems,skins,breakable]{tcolorbox}

% ---- Header/Footer ----
\usepackage{fancyhdr}
\renewcommand{\headrulewidth}{0.4pt}

% ---- Tables ----
\usepackage{booktabs}

% ---- Bibliography ----
\usepackage{natbib}
\bibliographystyle{plainnat}

% --------------------------------------------------------------------------
% Theorem environments
% --------------------------------------------------------------------------
\theoremstyle{plain}
\newtheorem{theorem}{Theorem}[section]
\newtheorem{lemma}[theorem]{Lemma}
\newtheorem{proposition}[theorem]{Proposition}
\newtheorem{corollary}[theorem]{Corollary}

\theoremstyle{definition}
\newtheorem{definition}[theorem]{Definition}
\newtheorem{example}[theorem]{Example}
\newtheorem{exercise}[theorem]{Exercise}
\newtheorem{assumption}[theorem]{Assumption}

\theoremstyle{remark}
\newtheorem{remark}[theorem]{Remark}

% --------------------------------------------------------------------------
% Colored boxes
% --------------------------------------------------------------------------
\newtcolorbox{keyresult}[1][]{%
  colback=green!5!white,
  colframe=green!50!black,
  fonttitle=\bfseries,
  title={Key Result},
  breakable,
  #1
}

\newtcolorbox{rlconnection}[1][]{%
  colback=blue!5!white,
  colframe=blue!50!black,
  fonttitle=\bfseries,
  title={RL/ML Connection},
  breakable,
  #1
}

\newtcolorbox{warning}[1][]{%
  colback=red!5!white,
  colframe=red!50!black,
  fonttitle=\bfseries,
  title={Warning},
  breakable,
  #1
}

% --------------------------------------------------------------------------
% Custom commands -- number sets
% --------------------------------------------------------------------------
\newcommand{\R}{\mathbb{R}}
\newcommand{\N}{\mathbb{N}}
\newcommand{\Q}{\mathbb{Q}}
\newcommand{\Z}{\mathbb{Z}}
\newcommand{\C}{\mathbb{C}}
\newcommand{\F}{\mathbb{F}}

% --------------------------------------------------------------------------
% Custom commands -- probability and expectation
% --------------------------------------------------------------------------
\newcommand{\E}{\mathbb{E}}
\newcommand{\Prob}{\mathbb{P}}
\newcommand{\Var}{\operatorname{Var}}
\newcommand{\Cov}{\operatorname{Cov}}
\newcommand{\indicator}{\mathbf{1}}

% --------------------------------------------------------------------------
% Custom commands -- convergence arrows
% --------------------------------------------------------------------------
\newcommand{\cas}{\xrightarrow{\text{a.s.}}}
\newcommand{\cprob}{\xrightarrow{\Prob}}
\newcommand{\clp}[1]{\xrightarrow{L^{#1}}}
\newcommand{\cdist}{\xrightarrow{d}}

% --------------------------------------------------------------------------
% Custom commands -- common operators
% --------------------------------------------------------------------------
\DeclareMathOperator{\dom}{dom}
\DeclareMathOperator{\epi}{epi}
\DeclareMathOperator{\conv}{conv}
\DeclareMathOperator{\relint}{relint}
\DeclareMathOperator{\aff}{aff}
\DeclareMathOperator{\cl}{cl}
\DeclareMathOperator{\interior}{int}
\DeclareMathOperator{\tr}{tr}
\DeclareMathOperator{\diag}{diag}
\DeclareMathOperator{\rank}{rank}
\DeclareMathOperator{\sgn}{sgn}
\DeclareMathOperator{\supp}{supp}
\DeclareMathOperator{\argmin}{arg\,min}
\DeclareMathOperator{\argmax}{arg\,max}
\DeclareMathOperator{\prox}{prox}
\DeclareMathOperator{\dist}{dist}
\DeclareMathOperator{\diam}{diam}
\DeclareMathOperator{\Span}{span}

% --------------------------------------------------------------------------
% Custom commands -- linear algebra operators
% --------------------------------------------------------------------------
\DeclareMathOperator{\nullity}{nullity}
\DeclareMathOperator{\range}{range}
\DeclareMathOperator{\nullspace}{null}
\DeclareMathOperator{\proj}{proj}

% --------------------------------------------------------------------------
% Custom commands -- measure theory
% --------------------------------------------------------------------------
\newcommand{\powerset}{\mathcal{P}}
\newcommand{\Borel}{\mathcal{B}}
 (from subdomain directories)

% ---- Encoding and fonts ----
\usepackage[utf8]{inputenc}
\usepackage[T1]{fontenc}

% ---- Mathematics ----
\usepackage{amsmath,amssymb,amsthm}
\usepackage{mathtools}

% ---- Page geometry ----
\usepackage[margin=1in,headheight=14pt]{geometry}

% ---- Hyperlinks ----
\usepackage[colorlinks=true,linkcolor=red,citecolor=blue,urlcolor=blue]{hyperref}

% ---- Lists ----
\usepackage{enumitem}

% ---- Colors and boxes ----
\usepackage{xcolor}
\usepackage[theorems,skins,breakable]{tcolorbox}

% ---- Header/Footer ----
\usepackage{fancyhdr}
\renewcommand{\headrulewidth}{0.4pt}

% ---- Tables ----
\usepackage{booktabs}

% ---- Bibliography ----
\usepackage{natbib}
\bibliographystyle{plainnat}

% --------------------------------------------------------------------------
% Theorem environments
% --------------------------------------------------------------------------
\theoremstyle{plain}
\newtheorem{theorem}{Theorem}[section]
\newtheorem{lemma}[theorem]{Lemma}
\newtheorem{proposition}[theorem]{Proposition}
\newtheorem{corollary}[theorem]{Corollary}

\theoremstyle{definition}
\newtheorem{definition}[theorem]{Definition}
\newtheorem{example}[theorem]{Example}
\newtheorem{exercise}[theorem]{Exercise}
\newtheorem{assumption}[theorem]{Assumption}

\theoremstyle{remark}
\newtheorem{remark}[theorem]{Remark}

% --------------------------------------------------------------------------
% Colored boxes
% --------------------------------------------------------------------------
\newtcolorbox{keyresult}[1][]{%
  colback=green!5!white,
  colframe=green!50!black,
  fonttitle=\bfseries,
  title={Key Result},
  breakable,
  #1
}

\newtcolorbox{rlconnection}[1][]{%
  colback=blue!5!white,
  colframe=blue!50!black,
  fonttitle=\bfseries,
  title={RL/ML Connection},
  breakable,
  #1
}

\newtcolorbox{warning}[1][]{%
  colback=red!5!white,
  colframe=red!50!black,
  fonttitle=\bfseries,
  title={Warning},
  breakable,
  #1
}

% --------------------------------------------------------------------------
% Custom commands -- number sets
% --------------------------------------------------------------------------
\newcommand{\R}{\mathbb{R}}
\newcommand{\N}{\mathbb{N}}
\newcommand{\Q}{\mathbb{Q}}
\newcommand{\Z}{\mathbb{Z}}
\newcommand{\C}{\mathbb{C}}
\newcommand{\F}{\mathbb{F}}

% --------------------------------------------------------------------------
% Custom commands -- probability and expectation
% --------------------------------------------------------------------------
\newcommand{\E}{\mathbb{E}}
\newcommand{\Prob}{\mathbb{P}}
\newcommand{\Var}{\operatorname{Var}}
\newcommand{\Cov}{\operatorname{Cov}}
\newcommand{\indicator}{\mathbf{1}}

% --------------------------------------------------------------------------
% Custom commands -- convergence arrows
% --------------------------------------------------------------------------
\newcommand{\cas}{\xrightarrow{\text{a.s.}}}
\newcommand{\cprob}{\xrightarrow{\Prob}}
\newcommand{\clp}[1]{\xrightarrow{L^{#1}}}
\newcommand{\cdist}{\xrightarrow{d}}

% --------------------------------------------------------------------------
% Custom commands -- common operators
% --------------------------------------------------------------------------
\DeclareMathOperator{\dom}{dom}
\DeclareMathOperator{\epi}{epi}
\DeclareMathOperator{\conv}{conv}
\DeclareMathOperator{\relint}{relint}
\DeclareMathOperator{\aff}{aff}
\DeclareMathOperator{\cl}{cl}
\DeclareMathOperator{\interior}{int}
\DeclareMathOperator{\tr}{tr}
\DeclareMathOperator{\diag}{diag}
\DeclareMathOperator{\rank}{rank}
\DeclareMathOperator{\sgn}{sgn}
\DeclareMathOperator{\supp}{supp}
\DeclareMathOperator{\argmin}{arg\,min}
\DeclareMathOperator{\argmax}{arg\,max}
\DeclareMathOperator{\prox}{prox}
\DeclareMathOperator{\dist}{dist}
\DeclareMathOperator{\diam}{diam}
\DeclareMathOperator{\Span}{span}

% --------------------------------------------------------------------------
% Custom commands -- linear algebra operators
% --------------------------------------------------------------------------
\DeclareMathOperator{\nullity}{nullity}
\DeclareMathOperator{\range}{range}
\DeclareMathOperator{\nullspace}{null}
\DeclareMathOperator{\proj}{proj}

% --------------------------------------------------------------------------
% Custom commands -- measure theory
% --------------------------------------------------------------------------
\newcommand{\powerset}{\mathcal{P}}
\newcommand{\Borel}{\mathcal{B}}
 (from subdomain directories)

% ---- Encoding and fonts ----
\usepackage[utf8]{inputenc}
\usepackage[T1]{fontenc}

% ---- Mathematics ----
\usepackage{amsmath,amssymb,amsthm}
\usepackage{mathtools}

% ---- Page geometry ----
\usepackage[margin=1in,headheight=14pt]{geometry}

% ---- Hyperlinks ----
\usepackage[colorlinks=true,linkcolor=red,citecolor=blue,urlcolor=blue]{hyperref}

% ---- Lists ----
\usepackage{enumitem}

% ---- Colors and boxes ----
\usepackage{xcolor}
\usepackage[theorems,skins,breakable]{tcolorbox}

% ---- Header/Footer ----
\usepackage{fancyhdr}
\renewcommand{\headrulewidth}{0.4pt}

% ---- Tables ----
\usepackage{booktabs}

% ---- Bibliography ----
\usepackage{natbib}
\bibliographystyle{plainnat}

% --------------------------------------------------------------------------
% Theorem environments
% --------------------------------------------------------------------------
\theoremstyle{plain}
\newtheorem{theorem}{Theorem}[section]
\newtheorem{lemma}[theorem]{Lemma}
\newtheorem{proposition}[theorem]{Proposition}
\newtheorem{corollary}[theorem]{Corollary}

\theoremstyle{definition}
\newtheorem{definition}[theorem]{Definition}
\newtheorem{example}[theorem]{Example}
\newtheorem{exercise}[theorem]{Exercise}
\newtheorem{assumption}[theorem]{Assumption}

\theoremstyle{remark}
\newtheorem{remark}[theorem]{Remark}

% --------------------------------------------------------------------------
% Colored boxes
% --------------------------------------------------------------------------
\newtcolorbox{keyresult}[1][]{%
  colback=green!5!white,
  colframe=green!50!black,
  fonttitle=\bfseries,
  title={Key Result},
  breakable,
  #1
}

\newtcolorbox{rlconnection}[1][]{%
  colback=blue!5!white,
  colframe=blue!50!black,
  fonttitle=\bfseries,
  title={RL/ML Connection},
  breakable,
  #1
}

\newtcolorbox{warning}[1][]{%
  colback=red!5!white,
  colframe=red!50!black,
  fonttitle=\bfseries,
  title={Warning},
  breakable,
  #1
}

% --------------------------------------------------------------------------
% Custom commands -- number sets
% --------------------------------------------------------------------------
\newcommand{\R}{\mathbb{R}}
\newcommand{\N}{\mathbb{N}}
\newcommand{\Q}{\mathbb{Q}}
\newcommand{\Z}{\mathbb{Z}}
\newcommand{\C}{\mathbb{C}}
\newcommand{\F}{\mathbb{F}}

% --------------------------------------------------------------------------
% Custom commands -- probability and expectation
% --------------------------------------------------------------------------
\newcommand{\E}{\mathbb{E}}
\newcommand{\Prob}{\mathbb{P}}
\newcommand{\Var}{\operatorname{Var}}
\newcommand{\Cov}{\operatorname{Cov}}
\newcommand{\indicator}{\mathbf{1}}

% --------------------------------------------------------------------------
% Custom commands -- convergence arrows
% --------------------------------------------------------------------------
\newcommand{\cas}{\xrightarrow{\text{a.s.}}}
\newcommand{\cprob}{\xrightarrow{\Prob}}
\newcommand{\clp}[1]{\xrightarrow{L^{#1}}}
\newcommand{\cdist}{\xrightarrow{d}}

% --------------------------------------------------------------------------
% Custom commands -- common operators
% --------------------------------------------------------------------------
\DeclareMathOperator{\dom}{dom}
\DeclareMathOperator{\epi}{epi}
\DeclareMathOperator{\conv}{conv}
\DeclareMathOperator{\relint}{relint}
\DeclareMathOperator{\aff}{aff}
\DeclareMathOperator{\cl}{cl}
\DeclareMathOperator{\interior}{int}
\DeclareMathOperator{\tr}{tr}
\DeclareMathOperator{\diag}{diag}
\DeclareMathOperator{\rank}{rank}
\DeclareMathOperator{\sgn}{sgn}
\DeclareMathOperator{\supp}{supp}
\DeclareMathOperator{\argmin}{arg\,min}
\DeclareMathOperator{\argmax}{arg\,max}
\DeclareMathOperator{\prox}{prox}
\DeclareMathOperator{\dist}{dist}
\DeclareMathOperator{\diam}{diam}
\DeclareMathOperator{\Span}{span}

% --------------------------------------------------------------------------
% Custom commands -- linear algebra operators
% --------------------------------------------------------------------------
\DeclareMathOperator{\nullity}{nullity}
\DeclareMathOperator{\range}{range}
\DeclareMathOperator{\nullspace}{null}
\DeclareMathOperator{\proj}{proj}

% --------------------------------------------------------------------------
% Custom commands -- measure theory
% --------------------------------------------------------------------------
\newcommand{\powerset}{\mathcal{P}}
\newcommand{\Borel}{\mathcal{B}}

\pagestyle{fancy}
\fancyhf{}
\fancyhead[L]{\textsc{Lesson 4: Function Sequences and Metric Spaces}}
\fancyhead[R]{\thepage}
\title{Lesson 4: Function Sequences and Metric Spaces}
\author{AI/ML Mathematics}
\date{}
\begin{document}
\maketitle
\tableofcontents
\newpage

%% ============================================================
%% SECTION 1: MOTIVATION AND OVERVIEW
%% ============================================================
\section{Motivation and Overview}
\label{sec:motivation}

Consider a reinforcement learning agent that iteratively improves its
value function. Starting from an arbitrary bounded function $V_0$, the
agent computes $V_1 = TV_0$, $V_2 = TV_1$, and so on, where $T$ is the
Bellman optimality operator. Three questions immediately arise: (1) Does
the sequence of functions $\{V_n\}$ converge? (2) If so, what kind of
convergence do we get---pointwise for each state, or something stronger?
(3) Is the limit guaranteed to be the unique fixed point of $T$? These
questions cannot be answered with sequences of real numbers alone; we
need a theory of \emph{sequences of functions} and, more fundamentally,
a notion of distance between functions. The framework of \emph{metric
spaces} provides exactly this: by equipping the set of bounded functions
with the supremum metric, we can treat value functions as points in a
complete metric space and apply the \emph{Banach fixed-point theorem} to
establish both existence, uniqueness, and a geometric convergence rate.
The tools developed in this lesson will let us prove that value
iteration converges, quantify its rate, and understand precisely when
limits commute with integration and other operations.

\begin{rlconnection}[title={Why RL/ML needs function sequences and metric spaces}]
\begin{itemize}[nosep]
  \item \textbf{Value iteration.} The Bellman operator $T$ is a
    $\gamma$-contraction on the complete metric space of bounded value
    functions under the sup-norm. The Banach fixed-point theorem
    guarantees convergence of $V_n = T^n V_0$ to the unique optimal
    value function $V^*$ at rate $O(\gamma^n)$.
  \item \textbf{Policy evaluation.} Temporal difference methods generate
    sequences of value function estimates; uniform convergence ensures
    that continuity and integrability of the limit are preserved.
  \item \textbf{Fitted value iteration.} When value functions are
    approximated by neural networks, one works in function spaces
    equipped with metrics (e.g., $L^2$ or sup-norm), and completeness
    of these spaces is essential for convergence guarantees.
  \item \textbf{Monte Carlo integration.} The interchange of limit and
    integral, justified by uniform convergence, underpins the validity
    of approximating expected rewards by sample averages of convergent
    function sequences.
\end{itemize}
\end{rlconnection}

\paragraph{Prerequisites.}
Convergence of real sequences, the Cauchy criterion, supremum and
infimum (Lesson~1, Definitions~3.4--3.5 and Theorem~4.10). Continuity
($\varepsilon$-$\delta$ characterization), uniform continuity, and the
extreme value theorem (Lesson~2, Theorems~5.2 and~7.4). The Riemann
integral, its properties, and the fundamental theorem of calculus
(Lesson~3, Theorems~\mbox{7.6--9.2}).

\paragraph{Learning objectives.}
After completing this lesson, you will be able to:
\begin{itemize}[nosep]
  \item \textbf{Define} pointwise and uniform convergence of function
    sequences, and \textbf{explain} the key difference (dependence of
    $N$ on $x$).
  \item \textbf{Prove} that the uniform limit of continuous functions is
    continuous and that uniform convergence permits interchange of limit
    and Riemann integral.
  \item \textbf{Apply} the Weierstrass $M$-test to determine uniform
    convergence of series of functions.
  \item \textbf{Define} metric spaces, open and closed sets, Cauchy
    sequences, completeness, and sequential compactness.
  \item \textbf{State} and \textbf{prove} the Banach fixed-point
    theorem, including existence, uniqueness, and the convergence rate
    estimate.
  \item \textbf{Justify} why value iteration converges by identifying
    the Bellman operator as a contraction on a complete metric space.
\end{itemize}

%% ============================================================
%% SECTION 2: REFERENCES
%% ============================================================
\section{References}
\label{sec:references}

\begin{itemize}
  \item \citet[Chapter 7]{rudin1976} --- Rigorous treatment of sequences
    and series of functions, uniform convergence, equicontinuity, and
    the Stone--Weierstrass theorem.
  \item \citet[Chapter 6]{abbott2015} --- Accessible development of
    pointwise vs.\ uniform convergence with many examples and the
    interchange theorems.
  \item \citet[Chapter 2]{rudin1976} --- Metric space topology: open
    and closed sets, compactness, completeness, and the contraction
    mapping principle.
  \item \citet[Chapters 8--9]{bartle2011} --- Detailed treatment of
    function sequences, metric spaces, and the contraction mapping
    theorem with applications.
  \item \citet[Chapter 2]{pugh2015} --- Metric spaces with a geometric
    viewpoint, completeness, and the Banach fixed-point theorem.
\end{itemize}

%% ============================================================
%% SECTION 3: POINTWISE AND UNIFORM CONVERGENCE
%% ============================================================
\section{Pointwise and Uniform Convergence}
\label{sec:pointwise_uniform}

We begin with the two fundamental modes of convergence for sequences of
functions. The distinction between them---seemingly a minor quantifier
swap---has profound consequences for the regularity of the limit.

\begin{definition}[Pointwise Convergence]
\label{def:pointwise}
Let $A \subseteq \R$ and let $\{f_n\}$ be a sequence of functions
$f_n \colon A \to \R$. We say $\{f_n\}$ \emph{converges pointwise} to
$f \colon A \to \R$ if for every $x \in A$ and every $\varepsilon > 0$,
there exists $N \in \N$ (which may depend on \emph{both} $x$ and
$\varepsilon$) such that
\[
  n \ge N \implies |f_n(x) - f(x)| < \varepsilon.
\]
We write $f_n \to f$ pointwise on $A$.
\end{definition}

\begin{definition}[Uniform Convergence]
\label{def:uniform}
Let $A \subseteq \R$ and let $\{f_n\}$ be a sequence of functions
$f_n \colon A \to \R$. We say $\{f_n\}$ \emph{converges uniformly} to
$f \colon A \to \R$ if for every $\varepsilon > 0$, there exists
$N \in \N$ (depending \emph{only} on $\varepsilon$) such that
\[
  n \ge N \implies |f_n(x) - f(x)| < \varepsilon
  \quad \text{for all } x \in A.
\]
Equivalently, $\sup_{x \in A} |f_n(x) - f(x)| \to 0$ as
$n \to \infty$. We write $f_n \rightrightarrows f$ on $A$.
\end{definition}

\begin{remark}
\label{rem:uniform_implies_pointwise}
Uniform convergence implies pointwise convergence (take the same $N$
for every $x$), but the converse is false, as the next example shows.
\end{remark}

\begin{example}[Pointwise but Not Uniform Convergence]
\label{ex:x_n}
Define $f_n \colon [0, 1] \to \R$ by $f_n(x) = x^n$. For each fixed
$x \in [0, 1)$, $x^n \to 0$ as $n \to \infty$, and $f_n(1) = 1$ for
all $n$. Thus $f_n \to f$ pointwise, where
\[
  f(x) = \begin{cases} 0 & \text{if } 0 \le x < 1, \\ 1 & \text{if } x = 1. \end{cases}
\]
Each $f_n$ is continuous on $[0, 1]$, but the limit $f$ is
\emph{discontinuous} at $x = 1$. We claim the convergence is not
uniform. Indeed, for any $n$, choosing $x_n = (1/2)^{1/n} \in [0, 1)$
gives $f_n(x_n) = 1/2$, while $f(x_n) = 0$, so
\[
  \sup_{x \in [0,1]} |f_n(x) - f(x)| \ge |f_n(x_n) - f(x_n)|
  = \tfrac{1}{2} \not\to 0.
\]
This example shows that pointwise convergence can destroy continuity.
Uniform convergence prevents this, as we prove next.
\end{example}

\begin{example}[Uniform Convergence on a Restricted Domain]
\label{ex:x_n_restricted}
Consider the same functions $f_n(x) = x^n$, but now on $[0, a]$ for a
fixed $a \in (0, 1)$. The pointwise limit is $f \equiv 0$, and
\[
  \sup_{x \in [0, a]} |f_n(x) - 0| = a^n \to 0
\]
as $n \to \infty$. Hence $f_n \rightrightarrows 0$ on $[0, a]$. The
limit is continuous, consistent with the theorem below.
\end{example}

\begin{keyresult}[title={Uniform Limit of Continuous Functions}]
\begin{theorem}[Continuity of the Uniform Limit]
\label{thm:uniform_limit_continuous}
Let $\{f_n\}$ be a sequence of continuous functions
$f_n \colon A \to \R$ that converges uniformly to $f$ on $A$. Then $f$
is continuous on $A$.
\end{theorem}
\end{keyresult}

\begin{proof}
Fix $c \in A$ and let $\varepsilon > 0$. We need to find $\delta > 0$
such that $|x - c| < \delta$ implies $|f(x) - f(c)| < \varepsilon$.
We split the difference into three parts:
\[
  |f(x) - f(c)| \le |f(x) - f_n(x)| + |f_n(x) - f_n(c)|
  + |f_n(c) - f(c)|.
\]
\emph{Step 1.} Since $f_n \rightrightarrows f$, choose $N$ such that
$\sup_{t \in A} |f_n(t) - f(t)| < \varepsilon/3$ for all $n \ge N$.
Fix any such $n$. Then $|f(x) - f_n(x)| < \varepsilon/3$ and
$|f_n(c) - f(c)| < \varepsilon/3$.

\emph{Step 2.} Since $f_n$ (for this fixed $n$) is continuous at $c$,
choose $\delta > 0$ such that $|x - c| < \delta$ implies
$|f_n(x) - f_n(c)| < \varepsilon/3$.

Combining:
\[
  |f(x) - f(c)| < \frac{\varepsilon}{3} + \frac{\varepsilon}{3}
  + \frac{\varepsilon}{3} = \varepsilon. \qedhere
\]
\end{proof}

\begin{warning}[title={Pointwise convergence does not preserve continuity}]
Example~\ref{ex:x_n} shows that pointwise convergence of continuous
functions can produce a discontinuous limit. The culprit is that the
convergence rate $N(x, \varepsilon)$ may blow up as $x$ approaches
certain points (in that example, $x \to 1^-$). Uniform convergence
eliminates this by requiring a single $N(\varepsilon)$ that works
simultaneously for all $x$. This is the function-sequence analogue of
the distinction between continuity and uniform continuity from Lesson~2.
\end{warning}

The next theorem shows that uniform convergence also permits
interchanging limits and Riemann integrals.

\begin{theorem}[Interchange of Limit and Integral]
\label{thm:interchange_limit_integral}
Let $\{f_n\}$ be a sequence of Riemann integrable functions on $[a, b]$
that converges uniformly to $f$ on $[a, b]$. Then $f$ is Riemann
integrable on $[a, b]$ and
\[
  \lim_{n \to \infty} \int_a^b f_n(x)\,dx = \int_a^b f(x)\,dx.
\]
\end{theorem}

\begin{proof}
\emph{Step 1: $f$ is Riemann integrable.} Given $\varepsilon > 0$,
choose $N$ such that $\sup_{x \in [a,b]} |f_n(x) - f(x)| <
\varepsilon / (3(b - a))$ for all $n \ge N$. Fix such an $n$. Since
$f_n$ is Riemann integrable, by the integrability criterion (Lesson~3,
Theorem~7.6), there exists a partition $P$ with
$U(f_n, P) - L(f_n, P) < \varepsilon/3$.

On each subinterval $[x_{i-1}, x_i]$, for any $s, t \in [x_{i-1}, x_i]$:
\begin{align*}
  f(s) - f(t) &= \bigl(f(s) - f_n(s)\bigr) + \bigl(f_n(s) - f_n(t)\bigr)
  + \bigl(f_n(t) - f(t)\bigr).
\end{align*}
Taking the supremum over $s$ and infimum over $t$:
\[
  M_i(f) - m_i(f) \le \frac{2\varepsilon}{3(b-a)} + (M_i(f_n) - m_i(f_n)),
\]
where $M_i(g) = \sup_{[x_{i-1}, x_i]} g$ and
$m_i(g) = \inf_{[x_{i-1}, x_i]} g$. Therefore:
\begin{align*}
  U(f, P) - L(f, P)
  &= \sum_{i=1}^k (M_i(f) - m_i(f))(x_i - x_{i-1}) \\
  &\le \frac{2\varepsilon}{3(b-a)} \sum_{i=1}^k (x_i - x_{i-1})
  + \sum_{i=1}^k (M_i(f_n) - m_i(f_n))(x_i - x_{i-1}) \\
  &= \frac{2\varepsilon}{3} + U(f_n, P) - L(f_n, P)
  < \frac{2\varepsilon}{3} + \frac{\varepsilon}{3} = \varepsilon.
\end{align*}
By the integrability criterion, $f$ is Riemann integrable.

\emph{Step 2: The integrals converge.} For $n \ge N$:
\begin{align*}
  \left|\int_a^b f_n - \int_a^b f\right|
  &= \left|\int_a^b (f_n - f)\right|
  \le \int_a^b |f_n - f| \\
  &\le (b - a) \sup_{x \in [a,b]} |f_n(x) - f(x)|.
\end{align*}
Since $\sup_{x \in [a,b]} |f_n(x) - f(x)| \to 0$, the right side
tends to $0$, giving $\int_a^b f_n \to \int_a^b f$.
\end{proof}

\begin{example}[Interchange Fails Under Pointwise Convergence]
\label{ex:interchange_fails}
Define $g_n \colon [0, 1] \to \R$ by
\[
  g_n(x) = \begin{cases}
    n^2 x & \text{if } 0 \le x \le 1/n, \\
    2n - n^2 x & \text{if } 1/n < x \le 2/n, \\
    0 & \text{if } 2/n < x \le 1.
  \end{cases}
\]
Each $g_n$ is a ``tent'' function with height $n$ and base $[0, 2/n]$.
For every fixed $x > 0$, $g_n(x) = 0$ for all $n > 2/x$, so
$g_n \to 0$ pointwise. Also $g_n(0) = 0$ for all $n$. Hence
$g_n \to 0$ pointwise on $[0, 1]$. But
$\int_0^1 g_n = \frac{1}{2} \cdot \frac{2}{n} \cdot n = 1$ for all $n$.
Thus
\[
  \lim_{n \to \infty} \int_0^1 g_n = 1 \neq 0
  = \int_0^1 \lim_{n \to \infty} g_n.
\]
The convergence is not uniform ($\sup |g_n| = n \to \infty$), so
Theorem~\ref{thm:interchange_limit_integral} does not apply.
\end{example}

\begin{rlconnection}[title={Uniform convergence and expected reward computation}]
In reinforcement learning, the expected cumulative reward under a policy
$\pi$ is $J(\pi) = \E\bigl[\sum_{t=0}^\infty \gamma^t r_t\bigr]$. When
we approximate the value function $V^\pi$ by a sequence of iterates
$V_n$, we need the approximation error in the integral
$\int V_n\,d\mu$ to vanish. Uniform convergence
$\|V_n - V^\pi\|_\infty \to 0$ guarantees
$\int V_n\,d\mu \to \int V^\pi\,d\mu$ for any probability measure
$\mu$. In Phase~01 (Measure Theory), the dominated convergence theorem
provides a more general sufficient condition that relaxes uniform
convergence to pointwise convergence plus a dominating integrable
function.
\end{rlconnection}

We now turn to series of functions, where uniform convergence is
established by a powerful comparison test.

%% ============================================================
%% SECTION 4: THE WEIERSTRASS M-TEST
%% ============================================================
\section{The Weierstrass $M$-Test}
\label{sec:weierstrass}

A \emph{series of functions} $\sum_{n=1}^\infty g_n(x)$ converges
pointwise (resp.\ uniformly) on $A$ if the partial sums
$S_N(x) = \sum_{n=1}^N g_n(x)$ converge pointwise (resp.\ uniformly)
on $A$.

\begin{theorem}[Weierstrass $M$-Test]
\label{thm:weierstrass_m_test}
Let $\{g_n\}$ be a sequence of functions $g_n \colon A \to \R$. Suppose
there exist constants $M_n \ge 0$ such that:
\begin{enumerate}[nosep,label=(\roman*)]
  \item $|g_n(x)| \le M_n$ for all $x \in A$ and all $n \ge 1$, and
  \item $\sum_{n=1}^\infty M_n$ converges.
\end{enumerate}
Then $\sum_{n=1}^\infty g_n(x)$ converges uniformly (and absolutely)
on $A$.
\end{theorem}

\begin{proof}
Let $S_N(x) = \sum_{n=1}^N g_n(x)$. We show $\{S_N\}$ is uniformly
Cauchy: for $M > N$,
\[
  |S_M(x) - S_N(x)| = \left|\sum_{n=N+1}^M g_n(x)\right|
  \le \sum_{n=N+1}^M |g_n(x)| \le \sum_{n=N+1}^M M_n.
\]
Since $\sum M_n$ converges, its tail $\sum_{n=N+1}^\infty M_n \to 0$ as
$N \to \infty$ (Lesson~1, the Cauchy criterion for series). Given
$\varepsilon > 0$, choose $N_0$ such that
$\sum_{n=N_0+1}^\infty M_n < \varepsilon$. Then for all $M > N \ge N_0$
and \emph{all} $x \in A$:
\[
  |S_M(x) - S_N(x)| \le \sum_{n=N+1}^M M_n
  \le \sum_{n=N_0+1}^\infty M_n < \varepsilon.
\]
Letting $M \to \infty$, the partial sums converge pointwise to some
function $S(x) = \sum_{n=1}^\infty g_n(x)$, and the bound above gives
$|S(x) - S_N(x)| \le \varepsilon$ for all $x \in A$ and $N \ge N_0$.
This means $\sup_{x \in A} |S(x) - S_N(x)| \le \varepsilon$ for
$N \ge N_0$, i.e., $S_N \rightrightarrows S$.

Absolute convergence follows because
$\sum |g_n(x)| \le \sum M_n < \infty$ for every $x \in A$.
\end{proof}

\begin{example}[Uniform Convergence of a Power Series]
\label{ex:power_series}
Consider the series $\sum_{n=0}^\infty x^n / n!$ on $[-R, R]$ for any
fixed $R > 0$. Setting $M_n = R^n / n!$, we have $|x^n / n!| \le M_n$
for all $|x| \le R$, and $\sum M_n = e^R < \infty$. By the Weierstrass
$M$-test, the series converges uniformly on $[-R, R]$. Since each
partial sum is a polynomial (hence continuous), the uniform limit
$e^x = \sum_{n=0}^\infty x^n / n!$ is continuous on $[-R, R]$ by
Theorem~\ref{thm:uniform_limit_continuous}.
\end{example}

\begin{example}[The $p$-Series of Functions]
\label{ex:p_series}
Define $g_n(x) = \sin(nx) / n^2$ on $\R$. Since
$|g_n(x)| \le 1/n^2 = M_n$ and $\sum 1/n^2 = \pi^2/6 < \infty$, the
Weierstrass $M$-test gives uniform convergence of
$\sum_{n=1}^\infty \sin(nx)/n^2$ on all of $\R$. Each partial sum is
continuous, so the limit function is continuous.
\end{example}

\begin{example}[A Series Where the $M$-Test Fails]
\label{ex:m_test_fails}
The series $\sum_{n=1}^\infty (-1)^{n+1} x^n / n$ converges pointwise
on $(-1, 1]$ to $\ln(1 + x)$ (by the alternating series test). However,
$\sup_{x \in (-1, 1]} |x^n / n| = 1/n$, and $\sum 1/n$ diverges, so
the $M$-test does not apply on the full interval. On any $[-a, a]$ with
$0 < a < 1$, we have $|x^n/n| \le a^n/n$ and $\sum a^n/n < \infty$
(comparison with a geometric series), so uniform convergence holds on
$[-a, a]$.
\end{example}

\begin{keyresult}[title={Connection to dominated convergence}]
The Weierstrass $M$-test is the series analogue of the dominated
convergence theorem from Lebesgue integration theory (Phase~01,
Lesson~5). Both results use a ``dominating'' bound ($M_n$ here, an
integrable function $g$ there) to justify interchange of limit and
summation or integration. In Phase~01, the dominated convergence theorem
will generalize Theorem~\ref{thm:interchange_limit_integral} from
uniform to pointwise convergence, provided a dominating integrable
function exists.
\end{keyresult}

With the theory of function sequences in hand, we now generalize the
notion of convergence itself by moving from $\R$ to abstract metric
spaces.

%% ============================================================
%% SECTION 5: METRIC SPACES
%% ============================================================
\section{Metric Spaces}
\label{sec:metric_spaces}

\begin{definition}[Metric Space]
\label{def:metric_space}
A \emph{metric space} is a pair $(X, d)$ where $X$ is a nonempty set
and $d \colon X \times X \to [0, \infty)$ is a function (called a
\emph{metric} or \emph{distance function}) satisfying, for all
$x, y, z \in X$:
\begin{enumerate}[nosep,label=(\roman*)]
  \item \textbf{Identity of indiscernibles:} $d(x, y) = 0$ if and only
    if $x = y$.
  \item \textbf{Symmetry:} $d(x, y) = d(y, x)$.
  \item \textbf{Triangle inequality:} $d(x, z) \le d(x, y) + d(y, z)$.
\end{enumerate}
\end{definition}

\begin{example}[Euclidean Metric on $\R^n$]
\label{ex:euclidean_metric}
The space $\R^n$ with the Euclidean metric
$d(\mathbf{x}, \mathbf{y}) = \bigl(\sum_{i=1}^n (x_i - y_i)^2
\bigr)^{1/2}$ is a metric space. The triangle inequality is the
Cauchy--Schwarz inequality applied to $\R^n$. When $n = 1$, this
reduces to $d(x, y) = |x - y|$, the standard metric on $\R$.
\end{example}

\begin{example}[Discrete Metric]
\label{ex:discrete_metric}
On any nonempty set $X$, define
\[
  d(x, y) = \begin{cases} 0 & \text{if } x = y, \\ 1 & \text{if } x \neq y. \end{cases}
\]
Verification of the metric axioms is immediate. The triangle inequality
holds because $d(x, z) \le 1 \le d(x, y) + d(y, z)$ whenever
$x \neq z$ (at least one of $d(x, y), d(y, z)$ equals $1$). In this
metric, every subset is open (since the open ball $B(x, 1/2) = \{x\}$),
and every sequence converges if and only if it is eventually constant.
\end{example}

\begin{example}[{Supremum Metric on $C[a, b]$}]
\label{ex:sup_metric}
Let $C[a, b]$ denote the set of continuous functions
$f \colon [a, b] \to \R$. Define
\[
  d_\infty(f, g) = \sup_{x \in [a, b]} |f(x) - g(x)|
  = \|f - g\|_\infty.
\]
This is a metric (the supremum exists by the extreme value theorem,
since $|f - g|$ is continuous on the compact set $[a, b]$). Note that
$d_\infty(f_n, f) \to 0$ is precisely the statement that
$f_n \rightrightarrows f$ on $[a, b]$: uniform convergence \emph{is}
convergence in the sup-metric.
\end{example}

\begin{definition}[Open Ball, Open Set, Closed Set]
\label{def:open_closed}
Let $(X, d)$ be a metric space.
\begin{enumerate}[nosep,label=(\roman*)]
  \item The \emph{open ball} of radius $r > 0$ centered at $x \in X$ is
    $B(x, r) = \{y \in X : d(x, y) < r\}$.
  \item A set $U \subseteq X$ is \emph{open} if for every $x \in U$
    there exists $r > 0$ such that $B(x, r) \subseteq U$.
  \item A set $F \subseteq X$ is \emph{closed} if its complement
    $X \setminus F$ is open.
\end{enumerate}
\end{definition}

\begin{proposition}[Properties of Open and Closed Sets]
\label{prop:open_closed_props}
In a metric space $(X, d)$:
\begin{enumerate}[nosep,label=(\roman*)]
  \item $\emptyset$ and $X$ are both open and closed.
  \item Arbitrary unions of open sets are open.
  \item Finite intersections of open sets are open.
  \item A set $F$ is closed if and only if every convergent sequence
    $\{x_n\} \subseteq F$ has its limit in $F$.
\end{enumerate}
\end{proposition}

\begin{proof}
(i) $\emptyset$ is open vacuously; $X$ is open since $B(x, 1) \subseteq
X$ for all $x$.

(ii) Let $U = \bigcup_{\alpha} U_\alpha$ where each $U_\alpha$ is open.
If $x \in U$, then $x \in U_\alpha$ for some $\alpha$, so there exists
$r > 0$ with $B(x, r) \subseteq U_\alpha \subseteq U$.

(iii) Let $U = U_1 \cap \cdots \cap U_k$ where each $U_i$ is open. If
$x \in U$, then $x \in U_i$ for each $i$, so there exist $r_i > 0$
with $B(x, r_i) \subseteq U_i$. Setting $r = \min(r_1, \ldots, r_k)
> 0$, we get $B(x, r) \subseteq U_i$ for all $i$, so
$B(x, r) \subseteq U$.

(iv) ($\Rightarrow$) Suppose $F$ is closed and $\{x_n\} \subseteq F$
with $x_n \to x$. If $x \notin F$, then $x \in X \setminus F$ which
is open, so there exists $r > 0$ with $B(x, r) \subseteq X \setminus
F$. But $x_n \to x$ implies $x_n \in B(x, r)$ for large $n$,
contradicting $x_n \in F$.

($\Leftarrow$) Suppose every convergent sequence in $F$ has its limit
in $F$. If $X \setminus F$ is not open, there exists
$x \in X \setminus F$ such that for every $n \ge 1$, $B(x, 1/n)$
contains a point $x_n \in F$. Then $d(x_n, x) < 1/n \to 0$, so
$x_n \to x$, forcing $x \in F$ by hypothesis---a contradiction.
\end{proof}

\begin{definition}[Convergence in a Metric Space]
\label{def:convergence_metric}
A sequence $\{x_n\}$ in a metric space $(X, d)$ \emph{converges} to
$x \in X$ if $d(x_n, x) \to 0$ as $n \to \infty$. That is, for every
$\varepsilon > 0$ there exists $N \in \N$ with
$d(x_n, x) < \varepsilon$ for all $n \ge N$. We write $x_n \to x$.
\end{definition}

\begin{proposition}[Uniqueness of Limits]
\label{prop:unique_limit}
In a metric space, limits are unique: if $x_n \to x$ and $x_n \to y$,
then $x = y$.
\end{proposition}

\begin{proof}
By the triangle inequality: $d(x, y) \le d(x, x_n) + d(x_n, y) \to
0 + 0 = 0$. Hence $d(x, y) = 0$, so $x = y$.
\end{proof}

With the basic metric space framework established, we can now define
Cauchy sequences and completeness---the properties that make convergence
arguments work.

%% ============================================================
%% SECTION 6: COMPLETENESS AND COMPACTNESS
%% ============================================================
\section{Completeness and Compactness}
\label{sec:completeness}

\begin{definition}[Cauchy Sequence]
\label{def:cauchy_metric}
A sequence $\{x_n\}$ in a metric space $(X, d)$ is a \emph{Cauchy
sequence} if for every $\varepsilon > 0$ there exists $N \in \N$ such
that $d(x_m, x_n) < \varepsilon$ for all $m, n \ge N$.
\end{definition}

\begin{proposition}
\label{prop:convergent_is_cauchy}
Every convergent sequence in a metric space is Cauchy.
\end{proposition}

\begin{proof}
If $x_n \to x$, then for $\varepsilon > 0$, choose $N$ with
$d(x_n, x) < \varepsilon/2$ for $n \ge N$. For $m, n \ge N$:
$d(x_m, x_n) \le d(x_m, x) + d(x, x_n) < \varepsilon/2 +
\varepsilon/2 = \varepsilon$.
\end{proof}

\begin{definition}[Complete Metric Space]
\label{def:complete}
A metric space $(X, d)$ is \emph{complete} if every Cauchy sequence in
$X$ converges to a limit in $X$.
\end{definition}

\begin{example}[$\R$ Is Complete; $\Q$ Is Not]
\label{ex:R_complete}
The real line $(\R, |\cdot|)$ is complete---this is the content of the
Cauchy convergence criterion (Lesson~1, Theorem~4.10). In contrast,
$\Q$ is not complete: the sequence $x_1 = 1, x_{n+1} =
\frac{1}{2}(x_n + 2/x_n)$ is Cauchy in $\Q$ (Newton's method applied
to $f(x) = x^2 - 2$), but its limit $\sqrt{2}$ is not in $\Q$.
\end{example}

\begin{theorem}[{$C[a, b]$ with the Sup-Metric Is Complete}]
\label{thm:C_complete}
The metric space $(C[a, b], d_\infty)$ is complete.
\end{theorem}

\begin{proof}
Let $\{f_n\}$ be a Cauchy sequence in $(C[a, b], d_\infty)$. We must
show it converges to some $f \in C[a, b]$.

\emph{Step 1: Pointwise limit exists.} For each fixed $x \in [a, b]$,
$|f_m(x) - f_n(x)| \le d_\infty(f_m, f_n)$, so $\{f_n(x)\}$ is a
Cauchy sequence in $\R$. Since $\R$ is complete, $f_n(x)$ converges to
some value, which we call $f(x)$. This defines a function
$f \colon [a, b] \to \R$.

\emph{Step 2: Convergence is uniform.} Given $\varepsilon > 0$, choose
$N$ such that $d_\infty(f_m, f_n) < \varepsilon$ for all $m, n \ge N$.
For any fixed $x \in [a, b]$ and $n \ge N$, take $m \to \infty$ in
$|f_m(x) - f_n(x)| \le \varepsilon$ (using continuity of
$|\cdot|$) to get $|f(x) - f_n(x)| \le \varepsilon$. Since this holds
for every $x$, $\sup_{x \in [a,b]} |f(x) - f_n(x)| \le \varepsilon$
for all $n \ge N$. Thus $f_n \rightrightarrows f$.

\emph{Step 3: $f$ is continuous.} Each $f_n$ is continuous and
$f_n \rightrightarrows f$, so $f$ is continuous by
Theorem~\ref{thm:uniform_limit_continuous}. Hence $f \in C[a, b]$ and
$d_\infty(f_n, f) \to 0$.
\end{proof}

\begin{rlconnection}[title={Completeness of function spaces in RL}]
The completeness of $(C[a,b], d_\infty)$ is a prototype for a central
pattern in reinforcement learning theory. The space of bounded functions
$B(S)$ on a state space $S$, equipped with the sup-norm
$\|V\|_\infty = \sup_{s \in S} |V(s)|$, is also complete (by the same
argument: pointwise limits of Cauchy sequences exist in $\R$, and the
convergence is uniform). This is the space in which value functions
live, and its completeness is what makes the contraction mapping
theorem applicable. In Phase~01 (Functional Analysis), we will study
Banach spaces---complete normed vector spaces---as the general framework
for these arguments.
\end{rlconnection}

\begin{definition}[Sequential Compactness]
\label{def:seq_compact}
A metric space $(X, d)$ is \emph{sequentially compact} if every
sequence in $X$ has a convergent subsequence (with limit in $X$).
\end{definition}

\begin{theorem}[Heine--Borel for $\R^n$]
\label{thm:heine_borel}
A subset $K \subseteq \R^n$ (with the Euclidean metric) is sequentially
compact if and only if it is closed and bounded.
\end{theorem}

\begin{proof}
($\Leftarrow$) Let $K \subseteq \R^n$ be closed and bounded, and let
$\{x_k\} \subseteq K$. Since $K$ is bounded, $K \subseteq [-M, M]^n$
for some $M > 0$. Write $x_k = (x_k^{(1)}, \ldots, x_k^{(n)})$. The
sequence $\{x_k^{(1)}\}$ is bounded, so by Bolzano--Weierstrass
(Lesson~1, Theorem~4.7), it has a convergent subsequence. Passing to
this subsequence, the second coordinates $\{x_k^{(2)}\}$ are bounded,
so they also have a convergent subsequence. Repeating $n$ times (each
time passing to a further subsequence), we obtain a subsequence
$\{x_{k_j}\}$ such that every coordinate converges. Hence
$x_{k_j} \to x$ for some $x \in \R^n$. Since $K$ is closed,
$x \in K$.

($\Rightarrow$) If $K$ is not bounded, choose $x_k \in K$ with
$\|x_k\| \ge k$; no subsequence can converge. If $K$ is not closed,
choose $x_k \in K$ with $x_k \to x \notin K$; every convergent
subsequence must also converge to $x \notin K$, contradicting
sequential compactness.
\end{proof}

\begin{example}
\label{ex:compact_interval}
The interval $[0, 1] \subset \R$ is sequentially compact (closed and
bounded). The open interval $(0, 1)$ is bounded but not closed: the
sequence $1/n \to 0 \notin (0,1)$ shows it is not sequentially compact.
The set $\R$ itself is closed but not bounded, hence not sequentially
compact: the sequence $n$ has no convergent subsequence.
\end{example}

We now have all the ingredients for the main result of this lesson.

%% ============================================================
%% SECTION 7: THE CONTRACTION MAPPING THEOREM
%% ============================================================
\section{The Contraction Mapping Theorem}
\label{sec:contraction}

\begin{definition}[Contraction Mapping]
\label{def:contraction}
Let $(X, d)$ be a metric space. A function $T \colon X \to X$ is a
\emph{contraction mapping} (or \emph{contraction}) if there exists a
constant $\gamma \in [0, 1)$ such that
\[
  d(T(x), T(y)) \le \gamma\, d(x, y) \quad \text{for all } x, y \in X.
\]
The constant $\gamma$ is called the \emph{contraction factor}.
\end{definition}

\begin{remark}
\label{rem:contraction_continuous}
Every contraction is (uniformly) continuous: given $\varepsilon > 0$,
taking $\delta = \varepsilon$ gives $d(T(x), T(y)) \le \gamma\,d(x,y)
< d(x,y) < \varepsilon$ whenever $d(x,y) < \delta$. In fact, every
contraction is Lipschitz continuous with Lipschitz constant $\gamma < 1$.
\end{remark}

\begin{keyresult}[title={Banach Fixed-Point Theorem}]
\begin{theorem}[Banach Fixed-Point Theorem (Contraction Mapping Theorem)]
\label{thm:banach}
Let $(X, d)$ be a complete metric space and let $T \colon X \to X$ be a
contraction with contraction factor $\gamma \in [0, 1)$. Then:
\begin{enumerate}[nosep,label=(\roman*)]
  \item \textbf{Existence.} $T$ has a fixed point: there exists
    $x^* \in X$ with $T(x^*) = x^*$.
  \item \textbf{Uniqueness.} The fixed point is unique.
  \item \textbf{Convergence rate.} For any starting point $x_0 \in X$,
    the iterates $x_n = T(x_{n-1}) = T^n(x_0)$ satisfy
    \[
      d(x_n, x^*) \le \frac{\gamma^n}{1 - \gamma}\, d(x_0, x_1).
    \]
\end{enumerate}
\end{theorem}
\end{keyresult}

\begin{proof}
Fix any $x_0 \in X$ and define $x_n = T(x_{n-1})$ for $n \ge 1$.

\textbf{Step 1: $\{x_n\}$ is Cauchy.} First, observe that
\[
  d(x_{n+1}, x_n) = d(T(x_n), T(x_{n-1}))
  \le \gamma\, d(x_n, x_{n-1}) \le \cdots \le \gamma^n\, d(x_1, x_0).
\]
For $m > n$, the triangle inequality gives:
\begin{align*}
  d(x_m, x_n)
  &\le \sum_{k=n}^{m-1} d(x_{k+1}, x_k)
  \le \sum_{k=n}^{m-1} \gamma^k\, d(x_1, x_0) \\
  &= \gamma^n \sum_{j=0}^{m-n-1} \gamma^j\, d(x_1, x_0)
  \le \gamma^n \sum_{j=0}^{\infty} \gamma^j\, d(x_1, x_0)
  = \frac{\gamma^n}{1 - \gamma}\, d(x_1, x_0).
\end{align*}
Since $\gamma \in [0, 1)$, $\gamma^n \to 0$, so the right side tends
to $0$ as $n \to \infty$. Hence $\{x_n\}$ is Cauchy.

\textbf{Step 2: Existence.} Since $(X, d)$ is complete, $x_n \to x^*$
for some $x^* \in X$. We show $T(x^*) = x^*$. Since $T$ is continuous
(Remark~\ref{rem:contraction_continuous}):
\[
  T(x^*) = T\left(\lim_{n \to \infty} x_n\right)
  = \lim_{n \to \infty} T(x_n) = \lim_{n \to \infty} x_{n+1} = x^*.
\]

\textbf{Step 3: Uniqueness.} Suppose $T(y^*) = y^*$ for some
$y^* \in X$. Then
\[
  d(x^*, y^*) = d(T(x^*), T(y^*)) \le \gamma\, d(x^*, y^*).
\]
Since $\gamma < 1$, this implies $(1 - \gamma)\, d(x^*, y^*) \le 0$,
so $d(x^*, y^*) = 0$, hence $x^* = y^*$.

\textbf{Step 4: Convergence rate.} From Step~1, for $m > n$:
$d(x_m, x_n) \le \frac{\gamma^n}{1 - \gamma}\, d(x_1, x_0)$.
Taking $m \to \infty$ and using continuity of the metric:
\[
  d(x^*, x_n) = \lim_{m \to \infty} d(x_m, x_n)
  \le \frac{\gamma^n}{1 - \gamma}\, d(x_1, x_0)
  = \frac{\gamma^n}{1 - \gamma}\, d(x_0, x_1). \qedhere
\]
\end{proof}

\begin{example}[Contraction on $\R$]
\label{ex:contraction_R}
Let $T \colon [0, \infty) \to [0, \infty)$ be defined by
$T(x) = \frac{1}{2}(x + 2/x)$ (Newton's method for $\sqrt{2}$). For
$x, y \ge 1$:
\[
  |T(x) - T(y)| = \left|\frac{1}{2}(x - y) - \frac{(x - y)}{xy}\right|
  = \frac{|x - y|}{2}\left|1 - \frac{2}{xy}\right|.
\]
On $[1, 2]$, $xy \ge 1$ and $2/(xy) \le 2$, so this analysis requires
care. Instead, note that $T$ maps $[\sqrt{2}, 2]$ into itself (since
$T(x) \ge \sqrt{2}$ by AM-GM and $T(x) \le x \le 2$ for
$x \ge \sqrt{2}$), and on $[\sqrt{2}, 2]$,
$|T'(x)| = \frac{1}{2}|1 - 2/x^2| \le 1/2$ by the MVT (Lesson~3,
Theorem~5.4). Hence $T$ is a contraction with $\gamma = 1/2$ on
$[\sqrt{2}, 2]$, and the unique fixed point is $x^* = \sqrt{2}$.
Starting from $x_0 = 2$: $x_1 = 3/2 = 1.5$, $x_2 = 17/12 \approx
1.4167$, converging rapidly to $\sqrt{2} \approx 1.4142$.
\end{example}

\begin{example}[Non-Example: $\gamma = 1$ Fails]
\label{ex:gamma_one}
The map $T \colon [0, \infty) \to [0, \infty)$ defined by $T(x) = x+1$
satisfies $|T(x) - T(y)| = |x - y|$ (so $\gamma = 1$, not strictly
less than $1$). It has no fixed point: $T(x) = x$ would require
$x + 1 = x$.

Similarly, the map $T(x) = x/(1+x)$ on $[0, \infty)$ satisfies
$|T(x) - T(y)| \le |x - y|$ (with $\gamma$ approaching $1$ as
$x, y \to 0$). It has a unique fixed point $x^* = 0$, but the
convergence rate is only $O(1/n)$, not geometric. The strict inequality
$\gamma < 1$ in the Banach theorem is essential for the geometric rate.
\end{example}

\begin{example}[Value Iteration as a Contraction]
\label{ex:value_iteration}
Consider a finite Markov decision process with state space
$S = \{1, \ldots, |S|\}$, action space $A$, rewards bounded by
$R_{\max}$, transition probabilities $p(s'|s,a)$, and discount factor
$\gamma \in (0, 1)$. The Bellman optimality operator
$T \colon \R^{|S|} \to \R^{|S|}$ is defined by
\[
  (TV)(s) = \max_{a \in A}\left[r(s, a) + \gamma \sum_{s' \in S}
  p(s'|s,a)\, V(s')\right].
\]
For any $V, W \in \R^{|S|}$ and any state $s$:
\begin{align*}
  |(TV)(s) - (TW)(s)|
  &= \left|\max_a\left[r(s,a) + \gamma \sum_{s'} p(s'|s,a) V(s')\right]
  - \max_a\left[r(s,a) + \gamma \sum_{s'} p(s'|s,a) W(s')\right]\right| \\
  &\le \max_a \gamma \sum_{s'} p(s'|s,a)\, |V(s') - W(s')| \\
  &\le \gamma \|V - W\|_\infty,
\end{align*}
where the first inequality uses $|\max_a f(a) - \max_a g(a)| \le
\max_a |f(a) - g(a)|$. Taking the supremum over $s$:
$\|TV - TW\|_\infty \le \gamma \|V - W\|_\infty$.

Since $(\R^{|S|}, \|\cdot\|_\infty)$ is a complete metric space and $T$
is a $\gamma$-contraction, the Banach fixed-point theorem gives:
\begin{enumerate}[nosep,label=(\roman*)]
  \item There exists a unique $V^* \in \R^{|S|}$ with $TV^* = V^*$
    (the Bellman optimality equation).
  \item For any initial $V_0$, the value iteration sequence
    $V_n = TV_{n-1}$ satisfies
    $\|V_n - V^*\|_\infty \le \frac{\gamma^n}{1 - \gamma}\|V_0 - V_1\|_\infty$.
\end{enumerate}
With $\gamma = 0.99$, reaching $\varepsilon$-accuracy requires about
$n \ge \frac{\log(\varepsilon(1-\gamma)/\|V_0 - V_1\|_\infty)}{\log \gamma}$
iterations.
\end{example}

%% ============================================================
%% SECTION 8: EXERCISES
%% ============================================================
\section{Exercises}
\label{sec:exercises}

\begin{exercise}[Sup-norm characterization of uniform convergence]
\textnormal{(\(*\))} \\
Let $f_n, f \colon [a,b] \to \R$ be bounded functions. Prove that
$f_n \rightrightarrows f$ on $[a,b]$ if and only if
$\|f_n - f\|_\infty \to 0$, where
$\|g\|_\infty = \sup_{x \in [a,b]} |g(x)|$.
\end{exercise}

\begin{exercise}[Weierstrass $M$-test application]
\textnormal{(\(*\))} \\
Show that the series $\sum_{n=1}^\infty \frac{\cos(nx)}{n^3}$ converges
uniformly on $\R$. Conclude that its sum is a continuous function.
\end{exercise}

\begin{exercise}[Completeness of $\ell^\infty$]
\textnormal{(\(**\))} \\
Let $\ell^\infty$ be the set of bounded real sequences
$\mathbf{x} = (x_1, x_2, \ldots)$ with the metric
$d(\mathbf{x}, \mathbf{y}) = \sup_{n \ge 1} |x_n - y_n|$. Prove that
$(\ell^\infty, d)$ is a complete metric space.

\textit{Hint:} Mimic the proof of Theorem~\ref{thm:C_complete}: show
that a Cauchy sequence in $\ell^\infty$ converges coordinatewise, the
convergence is uniform, and the limit is bounded.
\end{exercise}

\begin{exercise}[Failure of the contraction mapping theorem without completeness]
\textnormal{(\(**\))} \\
Let $X = (0, 1]$ with the standard metric $d(x, y) = |x - y|$ and
define $T \colon X \to X$ by $T(x) = x/2$.
\begin{enumerate}[nosep,label=(\alph*)]
  \item Show that $T$ is a contraction with factor $\gamma = 1/2$.
  \item Show that $T$ has no fixed point in $X$.
  \item Explain why this does not contradict the Banach fixed-point
    theorem.
\end{enumerate}
\end{exercise}

\begin{exercise}[Value iteration convergence rate]
\textnormal{(\(**\))} \\
Consider a two-state MDP with $S = \{s_1, s_2\}$, a single action,
rewards $r(s_1) = 1$, $r(s_2) = 0$, transition probabilities
$p(s_1 | s_1) = p(s_2 | s_2) = 1$ (absorbing states), and discount
factor $\gamma = 0.9$.
\begin{enumerate}[nosep,label=(\alph*)]
  \item Write the Bellman operator $T$ explicitly and verify that
    $V^* = (10, 0)$.
  \item Starting from $V_0 = (0, 0)$, compute $V_1, V_2, V_3$ and
    verify the bound
    $\|V_n - V^*\|_\infty \le \frac{\gamma^n}{1-\gamma}\|V_0 - V_1\|_\infty$.
  \item How many iterations are needed to guarantee
    $\|V_n - V^*\|_\infty < 0.01$?
\end{enumerate}
\end{exercise}

\begin{exercise}[Uniform convergence and differentiation]
\textnormal{(\(***\))} \\
The interchange of limit and derivative is more delicate than the
interchange of limit and integral. Consider $f_n(x) = \frac{\sin(nx)}{n}$
on $\R$.
\begin{enumerate}[nosep,label=(\alph*)]
  \item Show that $f_n \rightrightarrows 0$ on $\R$.
  \item Compute $f_n'(x)$ and show that $\{f_n'\}$ does not converge
    pointwise at $x = 0$.
  \item Conclude that uniform convergence of $\{f_n\}$ does not imply
    convergence of $\{f_n'\}$. Compare with
    Theorem~\ref{thm:interchange_limit_integral}: why is
    differentiation harder than integration?
\end{enumerate}
\textit{Hint:} For (c), note that differentiation amplifies oscillation
(multiplies by frequency $n$), while integration averages it out.
\end{exercise}

%% ============================================================
%% SECTION 9: SUMMARY
%% ============================================================
\section{Summary}
\label{sec:summary}

\begin{itemize}
  \item \textbf{Pointwise convergence} $f_n \to f$ requires
    $|f_n(x) - f(x)| < \varepsilon$ for each $x$ separately, with $N$
    depending on both $x$ and $\varepsilon$; it can destroy continuity
    of the limit (e.g., $x^n$ on $[0,1]$).

  \item \textbf{Uniform convergence} $f_n \rightrightarrows f$ requires
    a single $N(\varepsilon)$ that works for all $x$ simultaneously,
    equivalently $\|f_n - f\|_\infty \to 0$; it preserves continuity
    and permits interchange of limit and Riemann integral.

  \item The \textbf{Weierstrass $M$-test} establishes uniform
    convergence of a series $\sum g_n$ by bounding $|g_n(x)| \le M_n$
    with a convergent numerical series $\sum M_n$. It is the
    function-series analogue of the dominated convergence theorem.

  \item A \textbf{metric space} $(X, d)$ abstracts the notion of
    distance, enabling convergence, continuity, and completeness to be
    studied in general settings (function spaces, sequence spaces,
    etc.).

  \item \textbf{Completeness} (every Cauchy sequence converges) is the
    key property that distinguishes $\R$ from $\Q$ and makes fixed-point
    arguments work. The function space $(C[a,b], \|\cdot\|_\infty)$ is
    complete.

  \item \textbf{Sequential compactness} generalizes the
    Bolzano--Weierstrass theorem; in $\R^n$ (Heine--Borel), a set is
    sequentially compact if and only if it is closed and bounded.

  \item The \textbf{Banach fixed-point theorem} states that a
    contraction $T$ on a complete metric space has a unique fixed point
    $x^*$, and the iterates $T^n(x_0)$ converge to $x^*$ at a
    geometric rate $O(\gamma^n)$.

  \item The Bellman optimality operator is a $\gamma$-contraction on
    the complete space of bounded value functions, so \textbf{value
    iteration converges} to the unique optimal value function at rate
    $O(\gamma^n)$.

  \item Pointwise convergence can fail to commute with integration
    (tent function example) or differentiation ($\sin(nx)/n$); uniform
    convergence fixes the integration issue but not the differentiation
    one.
\end{itemize}

\bigskip

As the capstone lesson of the Real Analysis subdomain, we summarize the
full arc of Lessons~1--4 in the following table.

\begin{table}[h]
\centering
\caption{Real Analysis Subdomain: Lessons 1--4 Overview}
\label{tab:subdomain_summary}
\begin{tabular}{@{} c l l @{}}
\toprule
\textbf{Lesson} & \textbf{Core Topics} & \textbf{Key RL/ML Applications} \\
\midrule
1 & Real numbers, sequences, series &
  Convergence of iterative algorithms, \\
  & (completeness, Cauchy criterion, &
  existence of suprema/infima \\
  & Bolzano--Weierstrass) & for optimization \\
\midrule
2 & Topology, continuity, &
  Continuous dependence of policies \\
  & compactness on $\R$ &
  on parameters, extreme value \\
  & (EVT, IVT, uniform continuity) &
  theorem for bounded rewards \\
\midrule
3 & Differentiation, integration &
  Gradient descent (MVT, Taylor), \\
  & (chain rule, MVT, Taylor, &
  backpropagation (chain rule), \\
  & Riemann integral, FTC) &
  expected reward computation \\
\midrule
4 & Function sequences, metric spaces &
  Value iteration convergence \\
  & (uniform convergence, completeness, &
  (contraction mapping theorem), \\
  & Banach fixed-point theorem) &
  function space analysis \\
\bottomrule
\end{tabular}
\end{table}

%% ============================================================
%% SECTION 10: FURTHER READING
%% ============================================================
\section{Further Reading}
\label{sec:further_reading}

\begin{itemize}
  \item In \textbf{Phase~01} (Functional Analysis, Lessons~6--7), the
    metric space framework is extended to \emph{normed spaces} and
    \emph{Banach spaces} (complete normed spaces), providing the
    natural setting for the contraction mapping theorem in infinite
    dimensions. \emph{Hilbert spaces} add inner product structure,
    enabling orthogonal projections and the Riesz representation
    theorem.

  \item In \textbf{Phase~01} (Measure Theory, Lessons~4--5), the
    interchange of limit and integral is treated definitively via the
    \emph{monotone convergence theorem} and the \emph{dominated
    convergence theorem}, which generalize
    Theorem~\ref{thm:interchange_limit_integral} by replacing uniform
    convergence with weaker hypotheses.

  \item In \textbf{Phase~02} (Core RL, Lessons~13--14), the Banach
    fixed-point theorem is applied to the Bellman operators for both
    policy evaluation and optimal control, yielding convergence of
    value iteration and policy iteration with explicit rate bounds.

  \item \citet[Chapter 7]{rudin1976} continues with
    \emph{equicontinuity} and the \emph{Arzel\`{a}--Ascoli theorem},
    characterizing compact subsets of $C[a,b]$---essential for proving
    existence of optimal policies in continuous-state MDPs.

  \item \citet[Chapters 2--3]{pugh2015} develops metric space topology
    in greater depth, including the Baire category theorem and its
    applications to analysis.

  \item The Stone--Weierstrass theorem (\citet[Theorem 7.26]{rudin1976})
    shows that polynomials are dense in $C[a,b]$ under the sup-norm,
    providing a classical density result that foreshadows neural network
    universal approximation theorems (Phase~03).

  \item For historical context, Banach's original 1922 paper introduced
    the fixed-point theorem in the setting of complete metric spaces.
    The theorem was independently discovered by Caccioppoli in 1931.
    \citet[Chapter 9]{bartle2011} provides a textbook treatment with
    many applications.
\end{itemize}

\bibliography{real-analysis-refs}

\end{document}
